\subsection{A limitation of the fibre-dimension condition}

Let $Y$ be a closed definable subset of a closed bounded definable
space $X$.  An admissible mapping-cylinder description of a strict
tube $(V,Y,B)$ is a definable homeomorphism
\[
 M_\varphi=
 \bigl(B\times[0,1]\amalg Y\bigr)/((b,0)\sim\varphi(b))
 \longrightarrow V
\]
which is the identity on $Y$ and uses the prescribed identification
of $B\times\{1\}$ with the outer face $B$.  Here $V$ is a closed
subset of $X$, $V\setminus B$ is an open neighborhood of $Y$ in $X$,
and $B$ has the exterior collar required for a strict tube.
Write $q(b,t)$ for the associated map on $B\times[0,1]$.
In particular, its restriction at positive $t$ is a homeomorphism
onto $V\setminus Y$.

The obstruction comes from the junction of a disk and a transverse
segment. A finite attachment fibre would separate a connected
surface by finitely many points. We first justify the PL charts used
in that argument; a second proof below detects the same obstruction
by a winding circle and requires no manifold assumption on the
outer face.

\begin{lemma}[Compact marked charts and boundary preservation]
\label{lem:tube-marked-compact-charts}
Let $P$ be a compact polyhedron over an o-minimal real closed field.
A definable manifold chart at a point $p\in P$ gives a PL manifold
chart at $p$ in the existing PL structure on $P$.
For a boundary chart the conclusion respects its marked boundary.
Consequently a definable homeomorphism between compact definable
manifolds with boundary preserves their boundaries.
\end{lemma}
\begin{proof}
Subdivide so that $p$ is a vertex.  In the given chart choose the
inverse image $C$ of a small closed coordinate cube, with $p$ mapping
to its centre and with $C$ containing a relative neighborhood of $p$
in $P$.  Apply \cite[Theorem~2.2(1)--(2), author p.~6]{Shiota2010}
to the constant scalar function on $P$, marking $C$.
It gives a definable self-homeomorphism $\tau$ of $P$ preserving the
chosen triangulation and such that $C'=\tau^{-1}(C)$ is a
subpolyhedron.  Since $p$ is a vertex, $\tau(p)=p$, and $C'$ still
contains a neighborhood of $p$ in the original polyhedron.
The chart composed with $\tau$ identifies the compact marked pair
$(C',\{p\})$ definably with the closed cube and its centre.
Marked Hauptvermutung
\cite[Theorem~2.1, author p.~5]{Shiota2010}
provides a PL homeomorphism of these pairs.  It is a PL Euclidean
chart at the specified point $p$.

At a boundary point use a closed half-cube with $p$ in the relative
interior of its flat face, marking that face as well as $p$.
The same construction respects the boundary.  Thus a triangulation
of a compact definable manifold together with its designated boundary
is a PL manifold with that PL boundary.

Now let $h:M\to M'$ be a definable homeomorphism of two such compact
manifolds, and let $x\in M$.  In their compact polyhedral models,
subdivide to make $x$ and $h(x)$ vertices.  Apply marked
Hauptvermutung to the pairs $(M,\{x\})$ and $(M',\{h(x)\})$.
The resulting PL homeomorphism takes $x$ to $h(x)$ and preserves
intrinsic PL boundary.  Hence $x\in\partial M$ if and only if
$h(x)\in\partial M'$.  The argument does not assume beforehand
that $h$ preserves the boundary.
\end{proof}


\begin{theorem}\label{thm:tube-fibre-dimension-obstruction}
There is a compact nowhere dense semialgebraic subset $Y\subset S^3$,
with $Y\subset\overline{S^3\setminus Y}$, and a point $a\in Y$ such
that every admissible mapping-cylinder description of a strict tube
about $Y$ has
\[
 \dim\varphi^{-1}(a)>0,
 \qquad
 \dim_a S^3-\dim_aY-1=0.
\]
Thus the equality
\begin{equation}\label{eq:tube-fibre-dimension}
 \dim\varphi^{-1}(y)=\dim_yX-\dim_yY-1
\end{equation}
cannot always be imposed.  It holds for the link projection of a
locally trivial conical bundle over a pure-dimensional manifold.
Even when it holds, it does not imply contractibility of the derived
space of strict conical choices used in Theorem~\ref{result-12}.
\end{theorem}

\begin{proof}
In a stereographic chart of $S^3$, take
\[
 D=\{(x_1,x_2,0):x_1^2+x_2^2\leq1\},\qquad
 A=\{(0,0,z):0\leq z\leq1\},\qquad Y=D\cup A,
\]
and let $a=(0,0,0)$.  The stereographic inclusion gives a compact
semialgebraic subset of $S^3$.  Its interior is empty, and its local
dimension at $a$ is two.

Suppose $(V,Y,B)$ has an admissible description.  The positive product $B\times(0,1)$ is homeomorphic to an open
subset of $S^3$. Triangulate $B$ and take a compact product slab
$|B|\times[a,b]$ with $0<a<b<1$. At an interior point of each
vertical edge, the ambient definable three-manifold chart and
Lemma~\ref{lem:tube-marked-compact-charts} give a PL chart in
the product's existing PL structure. The link of the vertical
edge above a vertex $v$ of $B$ is, up to subdivision, the vertex
link of $v$ in $B$. It is a PL circle. Hence every vertex star
of $B$ is a disk, and $B$ is a closed definable surface without
boundary.

We next show that $B$ is connected.  Put
$r=(x_1^2+x_2^2)^{1/2}$.  For sufficiently small $\delta>0$, the
closure of
\[
 N=\{r<1+\delta,\ |z|<\delta\}
   \ \cup\ \{r<\delta,\ -\delta<z<1+\delta\}
\]
lies in $V\setminus B$ and in the chosen chart.  Such a $\delta$
exists because these neighborhoods shrink to the compact set $Y$.
The complement $N\setminus Y$ is path connected.  Here is an
explicit verification.  Below the disk, points can move radially
to the circle $r=1+\delta/2$ at height $-\delta/2$.  Above the disk,
points in the thin tube can first move off the axis, then descend
at a fixed positive radius smaller than $\delta$ to height
$\delta/2$.  Points in the enlarged disk neighborhood can move
within its upper half to that height.  They then move radially to
$r=1+\delta/2$.  This outer circle connects the upper and lower
parts through $z=0$ without meeting $D$.  Points initially on
$z=0$ outside $Y$ already have $r>1$.  These paths also cover the
axis above the tip of $A$, by first moving off that axis.

Every positive-radius ray of the mapping cylinder approaches a
point of $Y$, so it meets $N\setminus Y$.  Every point of
$V\setminus Y$ can therefore be joined to this connected set.
It follows that $V\setminus Y$, and hence $B$, is connected.

The map $q$ has compact, hence closed, image in $V$ and its image
contains $V\setminus Y$.  Since $Y$ is nowhere dense in $V$, this
image is all of $V$.  Positive-radius points map outside $Y$, so
$\varphi$ is onto.  The punctured set $Y\setminus\{a\}$ is the
disjoint union of the two nonempty relatively open and closed sets
$D\setminus\{a\}$ and $A\setminus\{a\}$.  Their inverse images
disconnect $B\setminus\varphi^{-1}(a)$.  A zero-dimensional
definable fibre is finite, whereas deleting finitely many points
from a connected surface does not disconnect it: paths can be
altered in finitely many disk charts to avoid those points.  This
contradiction proves $\dim\varphi^{-1}(a)>0$.

For the positive assertion, a local conical-bundle chart has the
form $U\times CL$, with $U$ open in a pure $m$-dimensional manifold
and the base given by $U$ times the cone vertex.  For nonempty $L$,
its local dimension along the base is $m+\dim L+1$, and the
attachment projection has fibre $L$.  This proves
\eqref{eq:tube-fibre-dimension}.  With the conventions
$C\varnothing=\{\mathrm{vertex}\}$ and
$\dim\varnothing=-1$, the same equality holds for the empty link.

Finally, take the point core of the cone on $T^6=(S^1)^6$.
For every strict neighborhood of its vertex, the punctured
neighborhood is an open seven-dimensional definable manifold and
is homeomorphic to $B\times(0,1)$. Thus $\dim B=6$.
Its attachment has the one vertex as target, so the sole fibre
is all of $B$ and satisfies \eqref{eq:tube-fibre-dimension}.
The numerical condition therefore includes the entire derived
component under consideration.  The derived choice component
specified in Theorem~\ref{result-12} has homotopy type
$BP_{\mathrm{def}}(T^6)$, whose fundamental group is nonzero by
Theorem~\ref{result-3}.  The numerical equality therefore does not
give categorical uniqueness in that precise derived model.
\end{proof}

\paragraph{Alternative proof of the dimension obstruction by winding.}
In a Euclidean chart of $S^3$, take
\[
 Y=\{(x,y,0):x^2+y^2\leq1\}
   \cup\{(0,0,z):0\leq z\leq1\},
 \qquad p=(0,0,0).
\]
This core is compact, semialgebraic, and nowhere dense.
At $p$ the proposed formula requires fibre dimension $3-2-1=0$.
Suppose an admissible tube
$q:B\times I\to V$ had a zero-dimensional fibre
$F=\varphi^{-1}(p)$.  The closed image of $q$ contains the dense
complement of $Y$ in $V$, so $q$ and $\varphi$ are onto and $F$ is
nonempty and finite.

Triangulate $(B,F)$ with the points of $F$ as vertices.
Small neighborhoods of the form $\{\lambda_f>c_f\}$, with
$c_f>1/2$, are contractible and pairwise disjoint; their straight-line
contractions to $f$ stay in their simplexes.
Choose them sufficiently small and choose $\delta>0$ so that, for
each $f\in F$,
\[
 q(N_f\times[0,\delta])\subset W,
\]
where $W$ is a small ambient ball about $p$, contained in $V$ and
of radius less than one in the displayed chart.
Continuity at the finitely many points $(f,0)$ supplies these choices.
The complement of
$\bigcup_f N_f\times[0,\delta)$ in $B\times I$ is closed and
bounded and its image under $q$ misses $p$.
Consequently there is a still smaller ball $W'$ about $p$ with
\[
 q^{-1}(W')\subset\bigcup_{f\in F}N_f\times[0,\delta).
\]
These are definable closed-image arguments over a general field.

Choose $a,h>0$ small enough that the horizontal circle
\[
 \gamma:S^1\longrightarrow W'\setminus Y,
 \qquad \gamma(u,v)=(au,av,h)
\]
lies in this smaller ball.
Its positive-radius inverse under $q$ is contained in one
$N_f\times(0,\delta)$, since $S^1$ is definably connected and the
sets in that finite union are disjoint and open.
This product is definably contractible, so applying $q$ gives a
nullhomotopy of $\gamma$ in $W\setminus Y$.

The plane section $W\cap\{z=0\}$ lies in the disk part of $Y$.
The connected image of the nullhomotopy must therefore remain in
$\{z>0\}$, where the normal segment removes the entire axis inside
$W$.  Angular projection
\[
 (x,y,z)\longmapsto
 \frac{(x,y)}{\sqrt{x^2+y^2}}
\]
is defined there and takes $\gamma$ to the identity of $S^1$.
This contradicts its nonzero degree, or equivalently its nontrivial
fundamental-group class.  In the arbitrary-field definable category
the same circle obstruction follows from
\cite[Corollary~3.4, author version p.~8]{BaroOtero2008v1}.
Thus every such fibre has positive dimension.
This argument imposes no manifold condition on $B$.

\paragraph{An explicit strict tube for this witness.}
There is also a direct finite construction showing that the example
has admissible strict tubes.
Choose a finite triangulation of $(S^3,Y)$, and barycentrically
subdivide so that the core is a full subcomplex $A$ of a finite
complex $L$.  For the semialgebraic witness the triangulation can
be chosen semialgebraic.  Work in this polyhedral model and write
\[
 w(x)=\sum_{v\in A}\lambda_v(x),\qquad
 V=\{w\geq1/2\},\qquad B=\{w=1/2\}.
\]
Fullness gives $w^{-1}(1)=|A|$, and $V$ contains the open
neighborhood $\{w>1/2\}$ of $|A|$.
For $b\in B$ put
\[
 a(b)=2\sum_{v\in A}\lambda_v(b)v,
 \qquad
 c(b)=2\sum_{v\notin A}\lambda_v(b)v.
\]
These are the normalized core and noncore parts of $b$, and
$b=(a(b)+c(b))/2$.
Define
\[
 q(b,t)=(1-t/2)a(b)+(t/2)c(b),\qquad 0\leq t\leq1,
 \qquad \varphi(b)=a(b).
\]
Both terms lie in faces of a common simplex, so this defines a
continuous semialgebraic map into $V$, with $q(b,1)=b$ and
$q(b,0)=\varphi(b)$.
For $x\in V\setminus|A|$, recover the two normalized parts by
\[
 a=\frac{\sum_{v\in A}\lambda_v(x)v}{w(x)},\qquad
 c=\frac{\sum_{v\notin A}\lambda_v(x)v}{1-w(x)}.
\]
The inverse of the positive-radius map is
\[
 x\longmapsto\left(\frac{a+c}{2},\,2(1-w(x))\right).
\]
It is continuous and semialgebraic.
Every point of the core lies in a simplex having a noncore coface:
the ambient complex is a triangulated three-manifold and the core
is nowhere dense.  Choosing a noncore vertex in such a coface
shows that $\varphi$ is onto.
The quotient identifies exactly those bottom points with the same
$\varphi$-value.  The finite polyhedral mapping-cylinder model and
closed-bounded definable compactness therefore identify its induced
map with a homeomorphism onto $V$, acting as the identity on the
specified copy of the core.
Finally the same displayed formula for $0<t<2$ has range
$\{0<w<1\}$ and the same inverse.
It gives a collar on both sides of $B$, including an exterior
collar for $t>1$.  The function $\rho=2(1-w)$ is zero exactly on
the core, and $V=\{\rho\leq1\}$, $B=\{\rho=1\}$.
Transporting this finite model back to $S^3$ proves strictness.


\subsection{Attachment fibres at a boundary core}

The real cell-like conclusion below is a boundary-core instance
of Lacher's relation between local connectivity of a mapping
cylinder and the neighborhood homotopy properties of its fibres
\cite[Theorem~2.2, pp.~652--653, and Corollary~2.4, p.~654]
{Lacher1970CellLikeII}.  We give the neighborhood argument explicitly
and then use definable triangulation to contract inside each fibre.
For an arbitrary o-minimal real closed field, all compactness,
homotopies, and manifold charts in the definable assertion have
their definable meanings.

\begin{theorem}\label{thm:boundary-attachment-fibres}
Let $V$ be a closed bounded definable manifold with boundary, let
$Y$ be a union of boundary components, and let $B$ be a closed
bounded definable set.  Suppose that a continuous definable map
\[
 q:B\times[0,1]\longrightarrow V
\]
satisfies $q(b,0)=\varphi(b)\in Y$ and restricts to a homeomorphism
$B\times(0,1]\longrightarrow V\setminus Y$.  Then $\varphi$ is
onto and each fibre $\varphi^{-1}(y)$ is definably contractible.

Over $\mathbb R$, suppose instead that $V$ is a compact topological
manifold with boundary and $B$ is a compact metric ANR, with the
same continuous-map hypotheses.  Then every fibre of $\varphi$ is
cell-like.  In the definable case over $\mathbb R$, these fibres
are contractible finite polyhedra up to definable homeomorphism.
\end{theorem}

\begin{proof}
The image of $q$ is closed and contains the dense set
$V\setminus Y$, so $q$ and $\varphi$ are onto.  

Fix $y\in Y$, put $F=\varphi^{-1}(y)$, and let $O$ be an open
neighborhood of $F$ in $B$, definable in the definable case.
The closed image $\varphi(B\setminus O)$ avoids $y$.  Choose an
open neighborhood $U$ of $y$ in $Y$ such that
$\varphi^{-1}(U)\subset O$.  For $0<\epsilon<1$, the set
\[
 N=U\ \cup\ q\bigl(\varphi^{-1}(U)\times[0,\epsilon)\bigr)
\]
is open in $V$. Indeed, its complement is the image under $q$ of
the closed bounded set
\[
 \bigl((B\setminus\varphi^{-1}(U))\times[0,1]\bigr)
 \ \cup\ \bigl(B\times[\epsilon,1]\bigr).
\]
That image is closed, also in the definable setting over an
arbitrary real closed field.  Choose a boundary half-ball $W$ about $y$ contained in $N$.
Because $Y$ is a union of boundary components, we may choose this
chart so that $W\cap Y$ is its whole flat boundary.  The set
$W\setminus Y$ is then homeomorphic to a convex open half-ball.
In particular, it has a contraction within itself, definable when
the data are definable.

There exists $0<\delta<\epsilon$ with
$q(F\times\{\delta\})\subset W\setminus Y$.  Indeed, the set
\[
 C=\{(b,t)\in F\times[0,1]:q(b,t)\notin W\}
\]
is compact and its projection to $[0,1]$ is closed and avoids zero.
Thus it is disjoint from an initial interval $[0,\delta]$ after
decreasing $\delta$.  This proof also applies to definably compact
sets and definable projections.

The open set
\[
 O_0=\{b\in B:q(b,\delta)\in W\}
\]
contains $F$.  Since $W\setminus Y\subset N\setminus Y$ and $q$
is injective at positive radius, $O_0\subset\varphi^{-1}(U)\subset O$.
Apply a contraction of $W\setminus Y$ to $q(b,\delta)$, and then
compose with
\[
 W\setminus Y\xrightarrow{\ q^{-1}\ }
 \varphi^{-1}(U)\times(0,\epsilon)
 \xrightarrow{\operatorname{pr}_B}\varphi^{-1}(U)\subset O.
\]
This is a nullhomotopy of the inclusion $O_0\hookrightarrow O$.
Its restriction is a contraction of $F$ inside $O$.

For the definable assertion, triangulate the compact pair $(B,F)$
using \cite[Theorem~2.2(1), p.~6]{Shiota2010}, and take a barycentric
subdivision.  Write the resulting pair as $(|K|,|L|)$, with $L$ a
full subcomplex of $K$.  If $a_v(x)$ denotes the global barycentric
coordinate of a vertex $v$ of $K$, put
\[
 O_L=\left\{x\in|K|:\sum_{v\in L}a_v(x)>0\right\},\qquad
 r(x)=\frac{\sum_{v\in L}a_v(x)v}{\sum_{v\in L}a_v(x)}.
\]
The vertices from $L$ having positive coordinates at $x$ span a
face in $L$, by fullness.  Hence $r$ is a continuous semialgebraic
retraction $O_L\longrightarrow|L|$.  The segment from $x$ to $r(x)$
stays in a simplex containing $x$.  Transport the retraction to
a definable neighborhood of $F$ in $B$, use the contraction of $F$
inside that neighborhood just constructed, and compose with the
retraction.  This contracts $F$ inside itself.

In the real topological case, the neighborhood construction proves
that $F\subset B$ has property $UV^\infty$.
Also, $B$ embeds as the compact slice $q(B\times\{1/2\})$ in $V$,
so it is finite dimensional.  A locally compact finite-dimensional
metric ANR is an ENR, and Lacher's criterion now makes $F$ cell-like
\cite[preliminary lemma, p.~718, and Theorem~1.1, p.~719]
{Lacher1969CellLike}.  This last conclusion requires no ANR
assumption on the fibre itself.
\end{proof}

\subsection{The numerical class of fixed-tube parametrizations}

Let $Y$ be a closed bounded definable manifold without boundary,
put $X=Y\times[0,2]$, and fix
\[
 V=Y\times[0,1],\qquad Y_0=Y\times\{0\},\qquad B=Y\times\{1\}.
\]
An admissible parametrization is a continuous definable map
$q:Y\times[0,1]\to V$ such that
\[
 q(y,1)=(y,1),\qquad q(y,0)=(\varphi(y),0),\qquad
 q|_{Y\times(0,1]}:Y\times(0,1]\xrightarrow{\sim}V\setminus Y_0.
\]
It induces the mapping-cylinder identification fixing the copy of
$Y_0$.  For the simplicial space of such maps, a simplex means a
jointly continuous definable family over the closed parameter
simplex, preserving that parameter.  Faces and degeneracies are
the usual affine pullbacks.  Write $P_{\mathrm{def}}(Y)$ for the
corresponding simplicial group of homeomorphisms of $Y\times[0,1]$
fixed pointwise on $Y\times\{0\}$.

\begin{theorem}\label{thm:numerical-tube-concordances}
The simplicial space of admissible parametrizations satisfying
\[
 \dim\varphi^{-1}(y)=\dim_yX-\dim_yY-1=0\qquad(y\in Y)
\]
is naturally isomorphic, by interval reversal, to
$P_{\mathrm{def}}(Y)$.  In particular, for $Y=T^6$ it has component
group
\[
 \bigoplus_{v\in\mathbb Z^6\setminus\{0\}}\mathbb F_2,
\]
and the numerical condition gives neither a connected nor a
contractible space of choices, even with $V$ and its outer-face
identification fixed.
\end{theorem}

\begin{proof}
By Theorem~\ref{thm:boundary-attachment-fibres}, every fibre of
$\varphi$ is nonempty and definably contractible.  Under the
dimension hypothesis each fibre is finite.  A nonempty finite
definably connected set is a point, so $\varphi$ is a continuous
definable bijection of compact sets, hence a homeomorphism.
It follows that $q$ is itself a continuous definable bijection
of compact sets, hence a homeomorphism of the cylinder fixing
its top pointwise.

Conversely, Lemma~\ref{lem:tube-marked-compact-charts} shows that
a definable homeomorphism of the cylinder preserves its whole
boundary over the given real closed field. If it fixes the top
pointwise, bijectivity prevents a bottom point from mapping to
the top. It therefore preserves the bottom setwise, including
when $Y$ is disconnected.  Its bottom restriction
is therefore a homeomorphism $\varphi$, and it supplies an
admissible parametrization with singleton fibres.  For a family
over a closed parameter simplex $\Delta$, the total map
\[
 \Delta\times Y\times[0,1]\longrightarrow
 \Delta\times Y\times[0,1],\qquad (s,y,t)\longmapsto(s,q_s(y,t))
\]
is a continuous definable bijection of compact sets.  Its inverse
is consequently jointly continuous and definable.  Thus the
identification holds for the stipulated families, as well as at
vertices.  Conjugation by $(y,t)\mapsto(y,1-t)$ changes the fixed
top to a fixed bottom and commutes with all affine parameter
pullbacks.  This is the required simplicial isomorphism.  The
torus calculation is Theorem~\ref{result-3}.  If $Y$ is empty,
both simplicial spaces are the one-point simplicial space.
\end{proof}

\paragraph{A noninjective attachment with a manifold mapping cylinder.}
The numerical hypothesis above has content.  Let the source circle
have cyclic vertices $b_0,b_1,b_2,b_3$ and the target circle have
cyclic vertices $c_0,c_1,c_2$.  Define a simplicial map by
\[
 b_0,b_1\longmapsto c_0,\qquad b_2\longmapsto c_1,\qquad
 b_3\longmapsto c_2.
\]
Its mapping cylinder is the annulus with the following seven
triangles:
\[
\begin{gathered}
 (b_0,b_1,c_0),\ (b_1,b_2,c_0),\ (b_2,c_1,c_0),\
 (b_2,b_3,c_1),\\
 (b_3,c_2,c_1),\ (b_3,b_0,c_2),\ (b_0,c_0,c_2).
\end{gathered}
\]
For each noncollapsed source edge the two triangles are the usual
triangulated square.  The collapsed edge has quotient parametrization
\[
 (u,t)\longmapsto (1-t)c_0+t\bigl((1-u)b_0+ub_1\bigr).
\]
This collapses exactly its bottom edge and is injective for $t>0$.
These parametrizations agree on the vertical edges, proving the
mapping-cylinder assertion.

For a direct planar verification of the annulus, assign coordinates
\[
\begin{aligned}
 b_0&=(-2,-2),&b_1&=(2,-2),&b_2&=(2,2),&b_3&=(-2,2),\\
 c_0&=(1,-1),&c_1&=(0,1),&c_2&=(-1,-1).
\end{aligned}
\]
The seven triangles have disjoint interiors, meet pairwise in
their indicated common faces, and meet the inner triangle only
in the indicated inner faces.  These assertions follow by the
linear half-plane inequalities of their edges.  Their areas, in
the order displayed, are
\[
 2,\quad2,\quad\tfrac52,\quad2,\quad\tfrac52,\quad2,\quad1.
\]
They sum to $14$, the area of the outer square minus the inner
triangle.  Their finite closed union is therefore exactly that
closed planar annulus: a missing interior point would leave a
nonempty open region of positive area, and both polygonal
boundaries are covered.  The outer four-cycle and inner
three-cycle are its two boundary components.  In particular,
the target circle is collared although the fibre over $c_0$ is
the nontrivial interval $[b_0,b_1]$.  This example is semialgebraic
over the rational numbers.  Its finite incidence inequalities
and quotient formulas remain valid over every real closed field.
Products give the same phenomenon with higher-dimensional
manifold boundary cores.

\subsection{Distinct components with unrestricted attachment fibres}

\begin{theorem}\label{thm:unrestricted-tube-components}
Let $Y=T^6=(S^1)^6$, put $V=Y\times[0,1]$, and fix its core
$Y_0=Y\times\{0\}$ and outer face $B=Y\times\{1\}$.
Let $\mathcal R$ be the space of all continuous maps $q:V\to V$
satisfying
\[
 q(y,1)=(y,1),\qquad q(Y_0)\subset Y_0,\qquad
 q|_{V\setminus Y_0}:V\setminus Y_0\xrightarrow{\sim}V\setminus Y_0,
\]
with the compact-open topology.  Its subspace of homeomorphisms
is $G=\operatorname{Homeo}(V\;\mathrm{rel}\;B)$.
There is a locally constant map
\[
 \lambda:\mathcal R\longrightarrow\pi_0G,
 \qquad\lambda(g)=[g]\quad(g\in G),
\]
where the target has the discrete topology.  Consequently
$\pi_0G\longrightarrow\pi_0\mathcal R$ is injective, and
$\mathcal R$ is not path connected.

Two fixed finite PL parametrizations witnessing this failure
also cannot be joined by a jointly continuous definable family
in any o-minimal expansion of $\mathbb R$, or by a semialgebraic
family over any real closed field.  Thus allowing arbitrary
attachment fibres does not give essential uniqueness for this
fixed strict tube in $Y\times[0,2]$.
\end{theorem}

\begin{proof}
We first record a metric observation.  Let $Z$ be a compact metric
space and $H$ a subgroup of $\operatorname{Homeo}(Z)$ for which
the identity path component contains an open neighborhood of the
identity.  The uniform metric
\[
 d_\infty(f,h)=\sup_{z\in Z}d(f(z),h(z))
\]
induces the compact-open topology on continuous self-maps and
on the homeomorphism group.  On homeomorphisms it is right
invariant.  In particular,
\begin{equation}\label{eq:tube-uniform-right-invariance}
 d_\infty(hg^{-1},\operatorname{id})=d_\infty(h,g).
\end{equation}
Choose $\epsilon>0$ such that every $h\in H$ with
$d_\infty(h,\operatorname{id})<\epsilon$ is in the identity
path component.  Equation~\eqref{eq:tube-uniform-right-invariance}
shows that any two elements of $H$ at distance less than
$\epsilon$ are in the same path component.

Let $\overline H$ denote its closure in $C(Z,Z)$.  Given
$f\in\overline H$, choose $h\in H$ with
$d_\infty(f,h)<\epsilon/4$ and set $\lambda_H(f)=[h]$.
Two such choices have distance less than $\epsilon/2$, so this
definition is independent of the choice.  If
$d_\infty(f,f')<\epsilon/4$, approximating homeomorphisms for
$f$ and $f'$ have distance less than $3\epsilon/4$, showing
$\lambda_H(f)=\lambda_H(f')$.  Thus $\lambda_H$ is locally
constant and extends the component map of $H$.

We apply this observation to $G$.  Edwards--Kirby local
contractibility applies to the homeomorphism group of the compact
manifold with boundary $V$
\cite[Corollary~1.1, p.~63; proof, p.~77]{EdwardsKirby1971}.
Here is the passage to the fixed-top subgroup.  The subgroup
$H\subset\operatorname{Homeo}(V)$ preserving both boundary
components is open.  If $h\in H$, write
$h(y,1)=(a_h(y),1)$ and define
\[
 R(h)=(a_h^{-1}\times\operatorname{id}_{[0,1]})\circ h.
\]
Restriction, inversion, and composition are continuous in the
compact-open group topology.  Hence $R:H\to G$ is continuous
and is the identity on $G$.  Take a sufficiently small local
contraction near the identity whose image lies in $H$ and
compose it with $R$.  It follows that a neighborhood of the
identity in $G$ contracts inside $G$.  In particular, $G$ has
the radius $\epsilon$ required above.

It remains to prove $\mathcal R\subset\overline G$.  Fix
$q\in\mathcal R$.  The real assertion of
Theorem~\ref{thm:boundary-attachment-fibres} shows that the
fibres of its bottom restriction are cell-like.  All other
fibres of $q$ are singletons.  The map $q$ is onto by compactness
and density, and is proper.  Moreover its restriction to the
whole boundary is a proper cell-like map onto the whole
boundary: it is the bottom attachment on $Y_0$ and the identity
on $B$.

Siebenmann's approximation theorem now applies to the compact
seven-manifolds with boundary
\cite[Theorem~A, p.~271]{Siebenmann1972Approximation}.
The theorem requires a proper cell-like map and a cell-like
boundary restriction.  Its condition~(a), which excludes
dimensions $3,4,5$, holds here.  Thus, for every $\eta>0$, there
is a homeomorphism $h:V\to V$ with
$d_\infty(h,q)<\eta$.

Use a product metric on $Y\times[0,1]$.  Choose $\eta$ smaller
than half the distance between the boundary components.
A homeomorphism preserves the manifold boundary and maps each
of its connected components onto a boundary component.  Since
$q|_B=\operatorname{id}_B$ and $d_\infty(h,q)<\eta$, no point
of $B$ can be sent to $Y_0$.  Hence $h(B)=B$ and $h(Y_0)=Y_0$.
The product metric and the same estimate on $B$ give
$d_\infty(a_h,\operatorname{id}_Y)<\eta$.  In addition,
\[
 d_\infty(a_h^{-1},\operatorname{id}_Y)
   =d_\infty(a_h,\operatorname{id}_Y)<\eta.
\]
Consequently $g=R(h)$ is in $G$ and
$d_\infty(g,q)<2\eta$.  This proves the required density
without making a continuous selection of approximations.
Restricting $\lambda_G$ to $\mathcal R$ gives $\lambda$.
Any path in $\mathcal R$ has constant $\lambda$, proving
injectivity on path components.

After interval reversal, $G$ is the ordinary topological
concordance group of $T^6$.  Theorem~\ref{result-3} gives
\[
 \pi_0G\cong
 \bigoplus_{v\in\mathbb Z^6\setminus\{0\}}\mathbb F_2.
\]
Choose a nonzero class represented by a finite PL concordance,
using the PL--TOP comparison in Theorem~\ref{result-4}.
By the finite coefficient model of
Theorem~\ref{thm:effective-countable-concordance-models}, we
may choose the representative $p$ over the field $F$ of real
algebraic numbers, on a fixed finite PL model of the torus.
Reversing the interval makes $p$ top-fixed.  The identity and
$p$ have different $\lambda$-values.  Both give admissible
mapping-cylinder identifications fixing the specified copy
of the core and the outer face.

A jointly continuous definable path over $\mathbb R$ is an
ordinary continuous path into $C(V,V)$, because $V$ is compact.
It therefore cannot join these two points of $\mathcal R$.
The same applies to a finite chain of definable one-simplices,
which can be concatenated using finitely many linear changes
of the parameter.

For the semialgebraic assertion, let $K$ be any real closed
field containing its canonical ordered copy of $F$.  Suppose
there were a semialgebraic family joining the two fixed
endpoints over $K$.  Its graph has a finite formula in the
language of ordered fields with a finite coefficient tuple
from $K$.  The conditions that it be a function, be jointly
continuous, have the specified endpoints, fix the top, send
the bottom into the bottom, and give a homeomorphism on the
complement of the bottom in every fibre are first-order
conditions on that tuple.  For the last condition one may
express bijectivity and continuity of the inverse graph by
the usual quantified $\varepsilon$--$\delta$ formulas.
All coefficients specifying the endpoints and the torus
belong to $F$.  Tarski's quantifier-elimination theorem makes
embeddings of real closed fields elementary
\cite[Theorem~3.2 and Corollary~3.3, pp.~2--3]{SwanTarski}.
Existential quantification of the finite tuple
and the elementary embedding $F\subset K$ give such a family
over $F$; the elementary embedding $F\subset\mathbb R$ then
gives a semialgebraic family over $\mathbb R$.  This contradicts
the preceding paragraph.  This transfer concerns semialgebraic
families over arbitrary real closed fields; the assertion for
general o-minimal families here is over $\mathbb R$.
\end{proof}

\begin{theorem}[Ordinary components of unrestricted rigidifications]
\label{thm:unrestricted-tube-component-bijection}
For the ordinary real six-torus, let
$V=T^6\times[0,1]$, let $B=T^6\times\{1\}$, and let
$G=\operatorname{Homeo}(V\;\mathrm{rel}\;B)$.
Let $\mathcal R$ have the compact-open topology and consist of the
continuous maps $q:V\to V$ that fix $B$ pointwise, map the bottom
into itself, and restrict to a homeomorphism off the bottom.
Every $q\in\mathcal R$ admits a continuous path
$\gamma:[0,1]\to\mathcal R$ with
\[
 \gamma(0)=q,\qquad \gamma(t)\in G\quad(0<t\leq1).
\]
Consequently inclusion induces a bijection
\[
 \pi_0G\xrightarrow{\ \cong\ }\pi_0\mathcal R.
\]
Under composition this is an isomorphism of component monoids;
in particular $\pi_0\mathcal R$ is the group identified with
$\pi_0G$.  The inverse bijection is induced by the locally constant
label of Theorem~\ref{thm:unrestricted-tube-components}.
\end{theorem}
\begin{proof}
The proof of Theorem~\ref{thm:unrestricted-tube-components} gives
$G\subset\mathcal R\subset\overline G$ for uniform convergence on
$V$, and gives a locally constant label
$\lambda:\mathcal R\to\pi_0G$ restricting to $g\mapsto[g]$.
Thus inclusion is already injective on path components.
We prove surjectivity with a controlled limiting isotopy.
This is the mechanism explicitly stated in the published Complement
to Theorem~A of
\cite[p.~272]{Siebenmann1972Approximation}.
We spell out the fixed-top construction.

For every $\epsilon>0$ there is $\delta>0$ such that any
$u\in G$ with $d_\infty(u,\mathrm{id})<\delta$ can be joined to
$\mathrm{id}$ inside the $\epsilon$-ball of $G$.
Indeed, Edwards--Kirby local contractibility
\cite[Corollary~1.1, pp.~63, 77]{EdwardsKirby1971}
provides contractions in arbitrarily small neighborhoods of the
identity in the manifold homeomorphism group.
Restrict to the open subgroup preserving the two boundary components
and compose its contraction with the continuous retraction
\[
 R(h)=(a_h^{-1}\times\mathrm{id}_{[0,1]})\circ h,
 \qquad a_h=h|_{T^6\times\{1\}}.
\]
Since $R$ fixes $G$ and the identity, continuity allows the
contraction to take values in any prescribed uniform neighborhood
in $G$.  This proves the claimed controlled path property.

For $n\geq1$ choose $\delta_n>0$ for
$\epsilon_n=2^{-n}$ and set
\[
 \eta_n=\min\{2^{-n},\delta_1/4,\ldots,\delta_n/4\}.
\]
For the specified $q\in\mathcal R$, density gives $g_n\in G$ with
$d_\infty(g_n,q)<\eta_n$.
Surjectivity of $g_n$ gives the exact identity
\[
 d_\infty(g_{n+1}g_n^{-1},\mathrm{id})
   =d_\infty(g_{n+1},g_n)<\delta_n/2.
\]
Choose a path $u_n$ from $g_{n+1}g_n^{-1}$ to the identity in the
$\epsilon_n$-ball, and put $\alpha_n(s)=u_n(s)g_n$.
It joins $g_{n+1}$ to $g_n$ and satisfies
\[
 d_\infty(\alpha_n(s),g_n)<\epsilon_n
 \qquad(0\leq s\leq1).
\]
On $[2^{-n},2^{-(n-1)}]$ define $\gamma$ by the affine
reparametrization of $\alpha_n$, with lower endpoint $g_{n+1}$
and upper endpoint $g_n$.  Define $\gamma(0)=q$.
The endpoints agree at every joining time.  On that interval,
\[
 d_\infty(\gamma(t),q)<\epsilon_n+\eta_n\longrightarrow0.
\]
Hence $\gamma$ is continuous at zero and everywhere else, with
$\gamma(1)=g_1$.  All positive-time maps are in $G$.
Compactness of $V$ makes this compact-open path equivalent to a
jointly continuous family $[0,1]\times V\to V$.
This proves surjectivity on path components and, with the label,
the stated bijection.

Composition preserves all defining conditions of $\mathcal R$
and is continuous for compact-open maps on compact $V$.
Inclusion is a monoid homomorphism, so the resulting bijection of
component monoids transports the group structure of $\pi_0G$.
\end{proof}

\paragraph{Scope.}
The paths just constructed are ordinary real paths obtained by
countably many successive isotopies.  No definable or semialgebraic
path-surjectivity, arbitrary-field analogue, higher weak equivalence,
or continuous choice of the path as $q$ varies is asserted.
Siebenmann's published Complement gives this limiting-isotopy
conclusion. Its subsequent announcement of nearly canonical
dependence on the map is not needed here.
